\documentclass[11pt]{article}
\usepackage[T1]{fontenc}
\usepackage{lmodern,amsmath,amssymb,amsthm,mathtools}
\usepackage[margin=0.86in]{geometry}
\usepackage{microtype,booktabs,longtable,array,xcolor,enumitem,fancyhdr}
\usepackage[hidelinks]{hyperref}
\usepackage{xurl}
\hypersetup{pdftitle={Reusing the CRT index in a Diophantine consistency equation},pdfauthor={AI-assisted research for Vladimir Reshetnikov}}
\setlength{\parskip}{4pt}
\setlength{\parindent}{0pt}
\setlist{nosep,leftmargin=1.6em}
\pagestyle{fancy}\fancyhf{}
\fancyhead[L]{\small Diophantine consistency: a reused CRT index}
\fancyhead[R]{\small 25 October 2026}\fancyfoot[C]{\thepage}
\setlength{\headheight}{14pt}
\newtheorem{theorem}{Theorem}[section]
\newtheorem{lemma}[theorem]{Lemma}
\newtheorem{proposition}[theorem]{Proposition}
\newtheorem{corollary}[theorem]{Corollary}
\theoremstyle{definition}\newtheorem{remark}[theorem]{Remark}
\newcommand{\Z}{\mathbb Z}\newcommand{\N}{\mathbb N}\newcommand{\Q}{\mathbb Q}
\newcommand{\Con}{\operatorname{Con}}
\newcommand{\lcm}{\operatorname{lcm}}
\newcommand{\LL}{\mathcal L}
% Generated from the literal equation.
\newcommand{\NV}{215}
\newcommand{\LastVar}{214}
\newcommand{\NR}{154}
\newcommand{\LastRow}{153}
\newcommand{\NT}{5660}
\newcommand{\Height}{91679}

\allowdisplaybreaks
\begin{document}
\begin{center}
{\LARGE\bfseries Reusing the CRT index in a\par Diophantine consistency equation}\\[9pt]
{\large An explicit quartic in \NV\ unrestricted integer variables}\\[10pt]
Prepared for Vladimir Reshetnikov and the ProveIt research program\\
AI-assisted mathematical research, 25 October 2026
\end{center}

\begin{abstract}
We construct an explicitly specified polynomial $P\in\Z[x_0,\ldots,x_{\LastVar}]$
of degree four such that ZFC proves that $P$ has no integer zero if and only
if ZFC is syntactically consistent. The polynomial is the sum of
\NR\ squares of quadratic polynomials, printed in the appendix and supplied
as exact coefficient data. Its collected expansion has \NT\ nonzero
monomials and coefficient height \Height.

The new arithmetic step eliminates a derived-formula-size test. The
independent CRT row-index representative already has a Pell pair and can
be chosen even and arbitrarily large in its residue class. It therefore
also supplies the semantic index bound used by the shared Pell certificate.
The removed size test is compensated by a larger label-weight parameter,
certified with a polynomially constructed exponent. A sharp elementary
product-separation lemma proves that this change preserves proof decoding.
Five exact quadratic projections and one additional affine projection
reduce the arithmetic source from 221 to 215 coordinates. Reversible
coordinate and row changes then reduce its coefficient height as well
as its number of monomials. The resulting 215-variable quartic is the
principal equation. The construction uses the numerical $C_{222}$ calculus directly.
The inherited logical calibration is identified explicitly; no new
end-to-end Lean or Rocq formalization is claimed.
\end{abstract}

\section{Result, comparison, and precise scope}

The principal equation is the following single ordinary Diophantine equation:
\begin{equation}\label{eq:main}
\boxed{\quad P(\boldsymbol x)=\sum_{j=0}^{\LastRow}R_j(\boldsymbol x)^2=0,
\qquad \boldsymbol x\in\Z^{\NV}.\quad}
\end{equation}
Every $R_j$ is an explicitly given quadratic polynomial with integer
coefficients. Appendix~\ref{app:literal} prints the entire list. The file
\nolinkurl{equations/principal/equation.json} defines the same list without
abbreviations, and its neighboring expansion contains every collected
coefficient of $P$. These files, the variable order, and the integer domain
are part of the definition. There is no uninstantiated universal polynomial,
machine index, existence-defined coefficient, or variable exponent in $P$.

\begin{theorem}\label{thm:main}
Using the identified inherited arithmetic and logical calibration,
the construction and proofs below establish
\[
\mathrm{ZFC}\vdash
\left(\neg\exists\boldsymbol x\in\Z^{\NV}\ P(\boldsymbol x)=0\right)
\longleftrightarrow\Con(\mathrm{ZFC}).
\]
Here consistency is the absence of a finite proof of falsity in the effective
ZFC presentation used by the repository's calibration.
\end{theorem}

The theorem is an equivalence, proved without assuming either consistency
or inconsistency. It is not a proof that the displayed polynomial has no
zeros. The traditional mathematical proof is separate from the finite
artifact checks. The latter authenticate coefficients, polynomial identities,
and reconstruction maps; they do not prove unbounded Pell or logical
calibration theorems by testing examples.

\begin{center}
\small
\begin{tabular}{@{}lrrrrr@{}}
\toprule Form & Variables & Squares & Degree & Monomials & Height\\\midrule
Pinned parent & 225 & 162 & 4 & 7,867 & 128,446\\
Principal & 215 & 154 & 4 & 5,660 & 91,679\\
Companion & 214 & 153 & 4 & 5,486 & $1.0996\times10^{12}$\\
Companion & 213 & 152 & 4 & 5,402 & $2.4179\times10^{24}$\\
\bottomrule
\end{tabular}
\end{center}
The companion heights are rounded upward here. Their exact values are
\[1\,099\,511\,627\,852\qquad\text{and}\qquad2\,417\,851\,639\,242\,452\,488\,945\,685.\]


The comparison is pinned to ProveIt commit
\nolinkurl{953fb4d82d4229f4efc5f9db4cc2bf38ae329bb4}.
The matched parent is the 225-variable pooled-power quartic in source
158~\cite{pooled}. Sources 151 and 157 supply the factorial-free binomial
CRT and affine beta constructions~\cite{factorialfree,linearcrt}; source
159 supplies the exact quadratic projection method~\cite{projection}.
The present result changes the choice and role of the semantic anchor and
removes the derived-label range, rather than merely reevaluating an old
polynomial. It also supplies a precise endpoint obtained by composing those
changes with exact algebraic reductions.

``Simplest'' is not a single ordered mathematical criterion. Variable
count, total degree, coefficient height, number of monomials, and circuit
size can improve in different directions. The repository contains
substantially lower-dimensional equations in other degree/coefficient
regimes. We claim a concrete improvement in the directly compiled quartic
family, not a global minimum or worldwide priority for the general idea.

\section{The retained proof-calculus architecture}

The numerical calculus has truth code $0$, falsity code $1$, modus ponens,
generalization, eight logical axiom schemata, and a guarded seed schema.
Its literal square-shell constructors are
\begin{align*}
 I(a,b)&=2((a+b)^2+a),&M(a,b)&=I(a,b)+1,\\
 A(a,b)&=2((a+b)^2+3a+4b+3),&E(a,b)&=A(a,b)+1.
\end{align*}
They encode implication, membership, universal quantification, and
equality. Put $N(a)=I(a,1)$ and $X(a,b)=N(A(a,N(b)))$.
The pairing maps $Q_0=(a+b)^2+a$ and
$Q_1=(a+b)^2+3a+4b+3$ enumerate complementary square shells;
parity supplies the four constructors. The inherited decoding theorem
transports finite refutations between these codes and the original
Cantor-coded $C_{222}$ calculus, fixing $0$ and $1$~\cite{shell}.

For clarity the retained logical codes are
\begin{align*}
 r_1&=X(a,E(a,b)),\\
 r_2&=I(A(a,I(b,c)),I(A(a,b),A(a,c))),\\
 r_3&=I(E(a,b),I(M(c,a),M(c,b))),\\
 r_4&=I(A(a,A(b,c)),A(b,A(a,c))),\\
 r_5&=I(E(c,a),I(M(c,b),M(a,b))),\\
 r_6&=I(X(a,A(a,b)),b),\\
 r_G&=I(E(c,a),I(E(c,b),A(M(a,b),E(a,b)))),\\
 r_8&=I(I(I(c,d),a),I(I(a,c),I(b,c))).
\end{align*}
Define
\begin{align*}
 \operatorname{Req}(b,k,m)
 &=A(b,I(m,N(A(b,I(m,N(M(b,k))))))),\\
 S(a,b,k)&=A(k,I(M(k,a),M(k,b))),\\
 T(a,b,k,m)&=N(A(a,I(S(a,b,k),\operatorname{Req}(b,k,m)))),\\
 \operatorname{Clos}(a,b,k,m)
 &=A(b,I(M(b,k),N(A(a,I(m,T(a,b,k,m)))))),\\
 \operatorname{nks}(a,b,c)&=A(c,I(\operatorname{Clos}(a,b,c,c),
                   N(\operatorname{Req}(b,c,c)))),\\
 \operatorname{RV}(a,b,c)&=I(M(a,b),A(c,M(a,b))),\\
 r_9&=N(I(\operatorname{RV}(a,b,c),\operatorname{nks}(a,b,c))).
\end{align*}
The last schema is allowed when $c\ne a,b$. Its arithmetic guard is
\[
 G=(c-a)^2(c-b)^2-(d+1)^2=0,\qquad a,b,c,d\in\N.
\]
Such a $d$ exists exactly when the two disequalities hold. The parameter
$d$ can separately be an ordinary formula parameter in the $r_8$ branch.
The alternatives do not require both interpretations simultaneously.

The proof transcript is encoded by positive weights $w(j)=1+Fj$.
A backward axiom step has $S'w(r)=S$, modus ponens has
$S'w(b)=S w(a)w(I(a,b))$, and generalization has
$S'w(A(a,b))=S w(b)$. Truth uses $S'=S$.
Let $D_r=S'w(r)-S$ and define the analogous MP and generalization
differences. The retained nonnegative local polynomial is
\[
 \LL=(S'-S)^2\!\prod_{r\in\mathcal R}D_r^2
       (D_{r_9}^2+G^2)D_{\rm MP}^2D_{\rm gen}^2+B_0^2+B_1^2.
\]
Here $\mathcal R=\{r_1,r_2,r_3,r_4,r_5,r_6,r_G,r_8\}$, and $B_0,B_1$ are the two beta-remainder equations. Thus $\LL=0$
means that one valid proof-rule branch and both beta equations hold.
The principal construction retains this local polynomial unchanged.

\section{Formula bounds without a range test}

The old arithmetic source separately required a derived sum of formula
labels to lie below the proof horizon in every local row. That required
its own numerator power, binomial divisibility test, and Pell index.
The replacement needs only a fixed, much more generous syntactic bound.

\begin{lemma}[Constructor-depth bound]\label{lem:depth}
For $t\ge16$, every formula label needed by a local rule with
$0\le a,b,c,d<t$ is smaller than $t^{2^{20}-1}$.
In particular it is at most $L=t^{2^{20}}$.
\end{lemma}
\begin{proof}
For a positive integer $s$, if integer inputs $0\le x,y<s$, then
$x,y\le s-1$ and each of
$I(x,y),M(x,y),A(x,y),E(x,y)$ is less than $8s^2$.
For the largest constructor,
$A(s-1,s-1)=8s^2-2s$, and $E$ adds one. If the inputs are
less than $t^e$, $e\ge1$, each output is consequently less than
$t^{2e+1}$. Starting with exponent one, a term of constructor depth
$h$ is bounded by $t^{2^{h+1}-1}$.

The exact syntax trees of the nine axiom codes have depths
\[
                 4,3,3,3,3,5,4,3,19.
\]
The MP and generalization labels have depth one; parameter labels have
depth zero. The depth certificate checks the trees against the pinned
literal constructor source, including all occurrences of $N$ and $X$.
The maximum depth nineteen proves the result. No asymptotic estimate or
coefficient rounding enters this argument.
\end{proof}

\begin{lemma}[Sharp product separation]\label{lem:separate}
Let $L,t,F$ be positive integers with $L\ge t\ge1$ and $F\ge L^t$.
Consider two finite lists of positive
integer labels at most $L$. Suppose the first has at most $t$ entries
and the second at most $2t+1$ entries. If their products of weights
$1+Fj$ agree, their label multisets agree.
\end{lemma}
\begin{proof}
First the lists have equal length. If their lengths differ, let the
smaller length be $r\le t$. The product for the smaller list is at most
$L^r(F+1)^r$, while the larger is at least $(F+1)^{r+1}$.
But $F+1>L^t\ge L^r$, a contradiction.

Write their common length as $m\le t$. For $1\le k\le m$, each
elementary symmetric coefficient $e_k$ of the labels satisfies
\[
 1\le e_k\le \binom{m}{k}L^k
       \le m^{m-k}L^k\le L^m\le F.
\]
The middle estimate uses $\binom{m}{k}=\binom{m}{m-k}\le m^{m-k}$.
Thus the corresponding coefficients of the difference of
$\prod(1+jX)$ have absolute value strictly less than $F$; the constant
coefficient is zero. If the polynomial difference were nonzero, divide
its value at $F$ by the lowest power with nonzero coefficient and reduce
modulo $F$. A nonzero integer of absolute value less than $F$ would be
divisible by $F$, impossible. The two polynomials coincide. Their linear
factors over $\Q$ recover every label with multiplicity. The empty-list
case is immediate. Labels zero, if present, contribute unit factors and
are discarded before this argument.
\end{proof}

The boundary $F=L^t$ is permitted. Bounding coefficient differences by
twice the largest coefficient would miss this useful fact: both coefficients
are positive, so their difference has absolute value at most $F-1$.

The new principal semantic choices are
\begin{equation}\label{eq:growth}
 t=128+4z^2,\qquad T_2=t^2,\qquad K=T_2^2=t^4,
 \qquad F=t^K.
\end{equation}
Since $t^3\ge128^3=2^{21}>2^{20}$, Lemma~\ref{lem:separate}
applies with $L=t^{2^{20}}$. Only the first three relations in
\eqref{eq:growth} are written as direct polynomial equations. The last
is certified by the retained Pell power mechanism. The exponent $K$
requires two quadratic wires, rather than a large coefficient multiplying
$t$. This is why the principal retains small coefficients.

\section{The factorial-free arithmetic front}

Use the following semantic quantities:
\begin{align}
 B&=\binom{F+t}{t},& N&=B^x,\label{eq:BN}\\
 y&=F(2v+1),& C&=2c_0,\label{eq:betaC}\\
 H&=2+4(C+a+b+c+d+S+S'+u+v_q),&q&=(BH)^x.\label{eq:radix}
\end{align}
The nine large CRT representatives are
$a,b,c,d,S,S',u,v_q,C$; their local residues encode one proof row.
The independent names $c_0$ and $v$ in \eqref{eq:betaC} are unrestricted
integer witnesses, not extra local schema parameters.

The endpoints and row-index congruence are
\begin{align}
 x&=F+1+(y+1)u_0,
 &x&=1+[1+(t+1)y]u_t,\label{eq:endpoints}\\
 B\eta&=C+F+1.\label{eq:selector}
\end{align}
For a local row with index $i=C\bmod p^{ex}$, the beta equations are
\[
 x=S+[1+(i+1)y]u,\qquad
 x=S'+[1+(i+2)y]v_q.
\]
The retained ranges are $a,b,c,d,C<t$, $S,S'\le y$, and $u,v_q<x$
on local decoded residues. There is no longer a range test on a derived
formula-label sum.

\subsection{Why one binomial modulus covers every row}

Here is the inherited selector theorem in the form used by this
construction~\cite{factorialfree}. Let $\Lambda_t=\lcm(1,\ldots,t)$.
If $R\ge\Lambda_t^2$ and $B=\binom{R+t}{t}$, then:
\begin{enumerate}
\item for every $1\le i\le t$, some full prime power $p^e\Vert B$
      divides $R+i$ and satisfies $p^e>t$;
\item every full prime power $p^e\Vert B$ divides $R+i$ for at least one
      such row $i$.
\end{enumerate}
Both directions matter. The first authenticates each row; the second
makes the reverse CRT construction compatible at every prime of $B$.

For completeness, the elementary proof is short. The valuation formula
shows $t\binom{t-1}{i-1}\mid\Lambda_t$. Since $R+i>\Lambda_t^2$,
some prime has $w=v_p(R+i)>2a$, where
$a=\lfloor\log_p t\rfloor$. Other factors $R+j$ have the same
$p$-valuation as $j-i$. Consequently
\[
 v_p(B)=w-v_p\left(t\binom{t-1}{i-1}\right)>a,
 \qquad v_p(B)\le w,
\]
which proves the first statement. For the second, let
$w=\max_i v_p(R+i)$. In Legendre's formula for $v_p(B)$ every summand
is zero or one, and all levels above $w$ vanish. Hence $v_p(B)\le w$;
a row attaining the maximum is compatible.

The standard elementary estimate $\Lambda_t\le4^{t-1}$ gives
$F\ge t^t\ge16^t>\Lambda_t^2$. Therefore $R=F$ is valid. Unlike
the parent, this construction has no equation requiring a variable to
equal a factorial.

\subsection{The range lemma and precision}

For a prime $p$, integers $A\ge0$ and $k,e\ge1$,
\begin{equation}\label{eq:rangelemma}
 p^e\mid\binom{A}{k}\quad\Longrightarrow\quad A\bmod p^e<k.
\end{equation}
If none of the $k$ consecutive factors $A-k+1,\ldots,A$ is divisible
by $p^e$, the valuation formula has at most $e-1$ nonzero levels, a
contradiction. A factor divisible by $p^e$ gives the claimed residue.
The convention $\binom{A}{k}=0$ for $A<k$ covers the boundary.

At a selector $p^e\Vert B$, the range for $C$ and
\eqref{eq:selector} give the exact local index $i-1$ in $0,\ldots,t-1$.
The other ranges recover the bounded local tuple modulo $p^{ex}$.
The retained local bound is
\begin{equation}\label{eq:localbound}
 0\le\LL<(2+t+x+y+F)^{2^{23}}<2^x\le p^{ex}.
\end{equation}
Its first inequality is the existing literal $C_{222}$ circuit
majorant~\cite{strictbudget}; changing the size of $F$ does not change
that polynomial bound. Once signs are recovered,
$x\ge2+(t+1)y$ and $y\ge F$ imply $2+t+x+y+F<2x$ and
$x\ge F+1>2^{32}$. If $2^n\le x<2^{n+1}$, $n\ge32$, then
$2^{23}(n+2)<2^n\le x$, proving the second strict inequality in
\eqref{eq:localbound}. Thus $N\mid\LL$ forces each selected local
value to be zero.

\section{A freely chosen CRT index as the Pell anchor}

The principal change is independent of the detailed syntax of $C_{222}$.
An integer used as a CRT representative can be enlarged by whole multiples
of its modulus without changing the finite transcript it represents.
We use that freedom to make $C$ the largest requested Pell index.

Write $X_j(A),Y_j(A)$ for the Pell polynomials, defined by
\[
 X_0=1,\quad Y_0=0,\qquad
 X_{j+1}=AX_j+(A^2-1)Y_j,\quad Y_{j+1}=X_j+AY_j.
\]
Signed indices have $X_{-j}=X_j$, $Y_{-j}=-Y_j$.
The square frame is $W=2+\sum\alpha^2$, where every actual power base,
index, and proposed output occurs among the $\alpha$. The retained odd
Pell frame proves $|A|>W^W$ and $|A-1|>4W^3$.

The fourteen nominal indices are
\begin{equation}\label{eq:indices}
 t,\ F+t,\ a,b,c,d,\ S,\ y+1,\ S',\ u,\ x,\ v_q,\ C,\ K.
\end{equation}
The former derived-label pair is reassigned to $K$; the existing $C$
pair becomes the anchor. With anchor coordinates $(U,Y)$, retain auxiliary
Pell pairs $(J_X,J_Y)$ and $(V_X,V_Y)$ and comparison center $D$ satisfying
\begin{align*}
 U^2-(A^2-1)Y^2&=1,&J_X^2-(A^2-1)J_Y^2&=1,\\
 V_X^2-(D^2-1)V_Y^2&=1,&J_Y&=Y^2(2e+1),\\
 D&=1+4J_Y(2d_0+1),&D&=A+J_Xf_0,\\
 V_X&=U+J_Xg_0,&V_Y&=C+4Yh_0.
\end{align*}
The squared bound is
\begin{equation}\label{eq:cap}
 Y^2=C^2+1+\sum_{i\ne\mathrm{anchor}}Y_i^2
             +r^2+s^2+w(w+1).
\end{equation}
All these variables are unrestricted integers. Multiplication wires turn
the displays into quadratic residuals; the final emitted data contain the
actual wires, not a higher-degree shorthand.

\begin{proposition}[Anchor substitution]\label{prop:anchor}
The retained signed anchor and pooled orientation proof applies with
nominal index $C$ in place of the derived formula-label index. No
monotonicity or syntactic interpretation of $C$ is required in soundness.
For completeness, $C$ can be chosen even and larger than all other indices
in \eqref{eq:indices} while preserving every required CRT residue.
\end{proposition}
\begin{proof}
Equation~\eqref{eq:cap} gives $|Y|>|C|$, so $Y\ne0$. Pell
classification and the retained squared-divisibility argument give an
absolute index $i>0$ with $|Y|=Y_i(|A|)$ and
$i=\pm C\pmod{2|Y|}$. Since $1\le i\le|Y|$ and $|C|<|Y|$,
one obtains $C\ne0$ and $i=|C|$. Every other classified index is
strictly smaller than $|C|$ by the same squared bound.

The pooled orientation row is, modulo its retained norm identities,
\[
 \sum_i\bigl((X_i-1)^2+(Y_i-k_i)^2\bigr)=(A-1)Q_0.
\]
Modulo $A-1$, write $s_i\in\{-1,1\}$ and $r_i$ for the small
abscissa and ordinate residues of each classified pair. The absolute
value of $r_i$ is its classified index, hence at most $|C|<W$.
Thus $|s_i-1|\le2$ and $|r_i-k_i|<2W$. For the fourteen pairs their
squared-error sum satisfies
\[
 0\le\sum_i\bigl((s_i-1)^2+(r_i-k_i)^2\bigr)
       <4\cdot14(W^2+1)<4W^3<|A-1|.
\]
It must vanish, identifying every pair with the canonical signed pair
at its nominal index. This includes the new $K$ pair and the $C$ anchor.
These are exactly the hypotheses and steps of the inherited signed-anchor
proof~\cite{anchor}; the name or origin of the nominal index never enters.

For the reverse construction, first fix $t,K,F,y,x,B,N$ and the
prescribed residue classes of the nine representatives. None of these
choices depends on the eventual representative $C$. Since $t$ is
divisible by four and $F=t^K$, $v_2(F)=K v_2(t)>\lfloor\log_2t\rfloor$.
Therefore $v_2(F+i)=v_2(i)$ for $1\le i\le t$, and $B$ and $N^2$
are odd. Each prescribed class modulo $N^2$ has cofinally many positive
even representatives, with step $2N^2$. Choose the other eight representatives
first, then choose $C$ in its class above all the other fixed indices.
This also preserves $B\mid C+F+1$, since $B\mid N^2$.
Only afterward are $H,q,W,A$ chosen. There is consequently no cyclic
size requirement.

For canonical pairs at a positive sufficiently large $A$, the strict index
inequality makes the slack in \eqref{eq:cap} nonnegative. Since $C$ is
even, $Y_C(A)$ is even and $X_C(A)$ odd. The inherited explicit auxiliary
choice at index $C Y_C(A)$ makes $J_Y/Y_C(A)^2$ odd; a two-modulus CRT
then constructs $D$ and the remaining comparison coordinates. Every
nonnegative slack has the ternary representation
$r^2+s^2+w(w+1)$. Thus all anchor witnesses exist.
\end{proof}

\section{Power certificates and the noncircular sign argument}

The polynomial retains sixteen power requests: three individual requests
for $F=t^K$, $N=B^x$, $q=(BH)^x$; ten at base $q+1$ for the binomial
numerators; and three at base $q$ for $q^t,q^{y+1},q^x$.
The square frame has thirty-five actual atoms, including duplicates where
the same integer has two roles. Its coefficient identity is independently
checked against every requested base, index, and output.

For a block at nonzero base $n$, let $Z_i=AY_i$ and let $m_i$ be the
proposed output. The inherited pooling mechanism uses just two rows:
\begin{align*}
 K_n-2\sum_iY_i(m_i-X_i)-(2A-n)L_n&=0,\\
 \sum_i((m_i-X_i+Z_i)^2-Y_i^2)-nK_n+L_n&=0.
\end{align*}
They imply divisibility of the sum of squared power errors by
$T_n=2An-n^2-1$. To handle signed indices, put
$b=\lfloor\sqrt{W-1}\rfloor$ and multiply by $n^{2b}$. Canonical
Pell congruences then give
\[
 T_n\mid E=\sum_i(m_i n^b-n^{b+k_i})^2,
 \qquad 0\le E\le4r b^{4b}<W^W<|T_n|.
\]
Here $r\le15$, every exponent $b+k_i$ is nonnegative, and
$\gcd(n,T_n)=1$. Consequently every output is exact; at $|n|\ge2$
the index is nonnegative, since a negative index would force an integer
times a positive power of $|n|$ to equal one. The individual requests
use the same canonical power argument. All displayed exponentials occur
in this proof, not as operations in the final polynomial.

It remains essential to justify the bases before interpreting these
requests. The order is as follows.
\begin{enumerate}
\item The known $t\ge128$ and $K=t^4$ give $F=t^K>0$, even.
      Since $C=2c_0$, equation~\eqref{eq:selector} has odd right side,
      so $B\ne0$. Independently, $H\equiv2\pmod4$ for every signed
      representative tuple. Thus $|BH|\ge2$.
\item The individual $BH$ request gives $x\ge0$ and $q=(BH)^x$.
      Equation $y=F(2v+1)$ gives $|y|\ge F$. If $x=0$, the last
      endpoint would make the integer $1+(t+1)y$, whose absolute value
      exceeds one, divide $-1$. Therefore $x\ge1$, $|q|\ge2$.
\item The $q$ block is now legitimate. Its $y+1$ index is nonnegative.
      A negative $y$ would be at most $-F<-1$, so $y\ge F$.
      The first endpoint excludes $u_0<0$, and the last excludes
      $u_t\le0$. Thus $x\ge F+1$ and $x\ge2+(t+1)y$.
\item Hence $|q+1|\ge2$. Its block recovers the nonnegative signs
      of all nine representatives and all their exact numerator powers.
      Now $H\ge2$. No earlier step needed these signs.
\item The first binomial digit authenticates $B>0$, as proved next.
      It follows that $q>0$ and $N=B^x>0$, completing the sign recovery.
\end{enumerate}

\subsection{Exact recovery of the digits}

Let $M=q^x$. Ten normalized remainder rows share
\begin{equation}\label{eq:remaindernorm}
 M^2=1+4\sum_{i=1}^{10}\rho_i^2+\alpha^2+\beta^2+\gamma(\gamma+1).
\end{equation}
For a denominator $D_i=q^{k_i}$, the ratio equation
$D_i a_i=M$ and normalized digit equation imply
$r_i=\rho_i/a_i\in\Z$ and the usual unscaled division equation.
The norm gives $|r_i|<|D_i|/2$, even before signs of the radix
are known. A true low binomial remainder with absolute value below
$|D_i|/32$ must equal $r_i$, by congruence and the strict interval width.

For the first digit, every coefficient of $(q+1)^{F+t}$ through degree
$t$ is at most $(F+t)^t\le2^{t(F-1)}$. On the other hand,
$|q|\ge2^x\ge2^{(t+1)F+2}$. The true low remainder and the
$t$-th coefficient are therefore smaller than their needed denominators
divided by $32$. The normalized first row gives
$B\equiv \binom{F+t}{t}\pmod q$. Also
$|B|\le|BH|<|q|/2$, since $x\ge3$; the positive binomial coefficient
is below $|q|/32$. This proves the exact identity
$B=\binom{F+t}{t}>0$.

For every other numerator with upper argument $a\ge0$ and requested
digit level $k\le x$,
$\binom{a}{j}\le\max(1,a)^x\le H^x$. Since $q=B^xH^x$ and
$B\ge2$, the same strict remainder bound follows. There is no
derived-label numerator needing a separate degree estimate here.
The ratios prove the necessary denominator divisibilities, and the
decoded digit equations yield the nine retained range predicates.

\section{Soundness and completeness of the new front}

\subsection{From an integer zero to a finite refutation}

The previous sections recover every sign, power, binomial digit, and
local range from an arbitrary integer zero. The selector theorem and
\eqref{eq:localbound} supply exactly $t$ valid local proof steps.
Their beta states join because the same $x,y$ determine the unique
remainders modulo $1+(i+1)y$.

The possible boundary $y=F$ must not be overlooked. Then the first beta
remainder is zero, whereas the last is one. Every valid local weight
equation propagates a zero state to zero, since all weights are positive.
This contradiction proves $y>F$. Thus the initial and terminal states
are respectively $F+1=w(1)$ and $1$.

Multiplying the local equations and cancelling the nonzero intermediate
states gives equality between the multiset of conclusion weights and
the initial falsity weight together with all premise weights. There are
at most $t$ nonunit conclusion factors and at most $2t+1$ nonunit premise
factors. All local labels are at most $L$ by Lemma~\ref{lem:depth}.
Lemma~\ref{lem:separate} gives equality of the corresponding label
multisets.

Finally, a balanced single-conclusion rule flow implies a finite
derivation of its distinguished target. To see this without assuming the
flow is acyclic, close the finitely many displayed rules under derivability,
starting with the valid axiom and truth conclusions. If falsity were outside
that closure, every rule with an outside conclusion would have at least one
outside premise. Summing ``outside conclusions minus outside premises''
over all occurrences would be at most zero. The recovered balance makes
the sum exactly one, contributed by the target falsity. Contradiction.
The finite closure computation itself extracts a derivation.

\subsection{From a finite refutation to an integer zero}

Start with a finite backward proof-tree expansion. Choose the cofinal
horizon $t=128+4z^2$ larger than its length and every local schema
parameter, including all freshness witnesses. There is no need to make
$t$ exceed the full expanded formula codes. Set $K=t^4$ and $F=t^K$;
form all finite pending-multiset products with weights $1+Fj$, and pad
the transcript with truth steps to length $t$.

Let $O$ be the odd part of $t!$. Choose an arbitrarily large positive odd
multiple $w$ of $O$ and put $y=Fw$. Since
$v_2(F)=K v_2(t)\ge2t>v_2(t!)$, one has $t!\mid y$.
Choose $w$ large enough that $y>F$ and every pending state is at most $y$.
The moduli $1+(i+1)y$, $0\le i\le t$, are pairwise coprime: a
common prime would divide both $y$ and one of these moduli after taking
an index difference at most $t$. CRT gives $x$ with the prescribed state
remainders; its residue class is cofinal, so all needed endpoint and
quotient inequalities can be met.

Now set $B=\binom{F+t}{t}$ and $N=B^x$. At every full prime power
$p^e\Vert B$, select one compatible row and choose all nine
representatives to match that row modulo $p^{2ex}$. The generalized CRT
theorem justifies this construction at every prime, including primes not
used as the distinguished selector of any row. Polynomial evaluation
preserves the local congruences. For a range with local value $w<k$,
the falling factorial contains the factor $A-w$, divisible by $p^{2ex}$;
the denominator loses at most $v_p(k!)\le k-1<x\le ex$ powers.
Thus the required binomial coefficient is divisible by $p^{ex}$.

Choose the eight ordinary representatives first. Then choose $C$ even,
positive, and larger than every other nominal index, in its unchanged
class modulo the odd integer $N^2$. Define $H$ and $q$ from their actual
values, so enlargement of $C$ strengthens the radix bounds. The ten true
normalized remainders are each less than $M/32$. Their common slack in
\eqref{eq:remaindernorm} is nonnegative, because
$4\cdot10/32^2<1$. Every integer $n\ge0$ has the form
$r^2+s^2+w(w+1)$: apply the three-square theorem to $4n+1$,
whose representation has one odd and two even coordinates.

Choose the large positive Pell center from the square frame, all canonical
pairs, and the even-$C$ auxiliary witnesses of Proposition~\ref{prop:anchor}.
The canonical congruences fill the individual and pooled power quotients.
This constructs an integer zero of the complete 221-variable quadratic
source. The exact projections below carry it to the final equation.

\section{Exact quadratic projections and coefficient control}

The semantic source has 221 coordinates and 160 quadratic residuals.
Its polynomial already has degree four and 7,361 collected monomials,
with height 128,446. All remaining reductions are exact algebra over
$\Z$, with explicit maps in both directions.

Write a quadratic residual as $r_i(x)=q_i(x)+\ell_i(x)+c_i$.
If an integral direction $v$ lies in the common nullspace of all Hessians
of the $q_i$, then
\[
 r_i(x+sv)=r_i(x)+s\ell_i(v).
\]
A coordinate of $v$ equal to $\pm1$ permits an integral unimodular
change of coordinates using $v$ as a column. If some $\ell_i(v)=\pm1$,
that residual solves the new coordinate as a quadratic polynomial in
the others. Substitution into the remaining rows preserves degree at
most two because the new coordinate occurs linearly with a constant
coefficient, namely $\ell_i(v)$, in each remaining row.
The inverse map is integral. This gives an existential equivalence over
all integers, not merely a rational projection.

Five such projections reduce the source to 216 variables and 155 rows.
One more affine consequence gives 215 variables and 154 rows. In the
216-stage coordinates, its exact relation is
\[
       R_{18}-R_5=4+12x_8+16x_9+4x_{19}-x_{20}+x_{33}.
\]
Solving for $x_{33}$ and removing the implied row is integral and degree
preserving. This immediate chart has large coefficients, but that
intermediate cost is removed by the reversible reductions described next.
All these operations are replays of literal certificates; the inherited
projection checker reconstructs forward and inverse identities independently
of the discovery algorithm.

Finally, integral Pell coordinate changes, integer row additions, and
integer translations reduce support and coefficient height. If $U$ is a
unimodular row matrix, the common zero set of $R$ equals that of $UR$;
the sums of their squares need not be the same polynomial. The resulting
principal equation has 215 variables, 154 residuals, 5,660 collected
monomials, and coefficient height 91,679. Its residual coefficient height
is 192. Thus the extra affine projection does not require accepting the
large coefficient height of its first chart. All reported final monomial
counts and heights belong to the fully collected polynomial after these
changes.

\section{Companions and an explored alternative}

The package also provides 213- and 214-variable quartic companions,
obtained by choosing $K=2^{20}t$ and $t=16+4z^2$.
Lemma~\ref{lem:separate} still applies, including its boundary.
Substituting $K$ directly into the square frame gives the 213-variable
version with large coefficients. Retaining $K$ as one separate affine
coordinate gives the 214-variable version and reduces that cost.
The principal 215-variable version instead uses two quadratic wires for
$K=t^4$ and attains the lower coefficient height reported above.
The comparison table gives these alternatives under the same unrestricted
integer domain and total-degree conventions.

A separate logical experiment splits the old guarded seed bundle into
guarded membership vacuity and the one-parameter schema
\[
             N(\operatorname{nks}(0,1,c+2)).
\]
The new axioms are old consequences. Conversely, the inherited full-seed
completeness proof uses the seed only at roles $0,1$ with a private code
greater than one, while membership vacuity supplies the remaining logical
interface. The accompanying proof note explains this ordinary calibration.
Under square-shell recoding, the parameter $c+2$ must be transported by
recoding the private code and subtracting two; addition itself is not
preserved by the code bijection.

This experiment lowers the formula-depth budget but its independent
compiler produced 229 variables where the matched old source uses 225.
It loses some sharing among formula constructors. It is retained as a
proved research observation, not presented as a smaller principal equation.
The principal result uses the original calibrated $C_{222}$ calculus.

\section{Logical calibration and what was verified}

The semantic chain is
\[
\begin{gathered}
 \text{integer zero of the printed polynomial}\\
 \Longleftrightarrow\text{integer zero of the new arithmetic source}\\
 \Longleftrightarrow\text{finite square-shell }C_{222}\text{ refutation}\\
 \Longleftrightarrow\text{finite original }C_{222}\text{ refutation}\\
 \Longleftrightarrow\text{finite ZF refutation}
 \Longleftrightarrow\text{finite ZFC refutation}.
\end{gathered}
\]
The last two equivalences are the named repository interfaces
\nolinkurl{KEndpoint.c222_iff_zf_refutation} and
\nolinkurl{ConstructibleChoice.zf_refutation_iff_zfc_refutation}.
The latter removes Choice using the constructible-inner-model argument
and soundness/completeness. The former uses the seed-model lower chain
and the direct upper interpretation~\cite{calibration}.
These endpoint source files were inspected, but their entire transitive
Lean dependency closure was not rebuilt in this iteration.

The new contribution is an ordinary mathematical proof, formalizable in
ZFC, combined with exact finite computations. The archive's verification
driver checks the new source against independently reconstructed semantic
formulas, all sixteen power requests and their frame atoms, the signed
Pell offsets, the relocated anchor, the removed normalized remainders,
every exact projection, and all final coordinate and row changes. It also
collects every quartic coefficient and checks the printed appendix against
the machine-readable residuals. A rebuild runs in a fresh temporary
directory. The same substantive audits run with Python optimization enabled,
so correctness checks do not depend on removable \texttt{assert} statements.

The inherited unbounded Pell classification, squared-divisibility and
comparison-index lemmas, the local-polynomial degree ledger, and the
$C_{222}$/ZF/ZFC calibration remain mathematical dependencies. Their
locations and copied proof sources are identified in the provenance files.
Finite checking is not described as proving these dependencies anew.

\section{Further research questions}

\begin{enumerate}
\item Can the existing exponent set provide a sufficiently large weight
      parameter without either the extra $K$ index or the two $t^4$ wires?
      A useful candidate must preserve the acyclic order: $F$ precedes
      the beta code, which precedes the representative $C$ and radix.
\item Can the principal 215-variable form admit another integral projection
      while keeping its coefficients small? Its common Hessian radical is
      zero and its 154 quadratic row parts are linearly independent.
      Both properties persist under invertible affine coordinate changes
      and constant invertible row operations, so these two projection
      methods are exhausted. Other polynomial substitutions and different
      arithmetic encodings remain possible.
\item Can a logic-specific encoding realize the smaller fixed-role seed
      depth while recovering the arithmetic sharing lost by its two-schema
      split? The present negative measurement gives a concrete comparator.
\item Can the anchor-transport theorem be formulated once for a general
      family of arithmetic compilers, with a machine-checked side condition
      certifying an unbounded CRT representative class and an available
      canonical pair?
\item The best next formalization target is the new source-to-transcript
      equivalence and its exact certificate replay. This is more decisive
      than additional bounded tests of Pell identities. A future complete
      endpoint should compose it with a reproducible build of the inherited
      ZFC calibration.
\end{enumerate}

\appendix
\section{The literal principal equation}\label{app:literal}

All variables $x_0,\ldots,x_{\LastVar}$ range independently over $\Z$.
The equation is $\sum_{j=0}^{\LastRow}R_j^2=0$ with the following
complete residual list. A missing monomial has coefficient zero.
The displayed expressions are parsed back to integer polynomials by
the release checker. Coordinate names in construction-stage metadata
are historical; the exact inverse certificates give their semantic meaning.

% Full literal residual list; parsed independently by the release checker.
{\small
\setlength{\jot}{2pt}
\allowdisplaybreaks[0]
% BEGIN RESIDUAL 0
\begin{align*}
R_{0} ={}& x_{0} - x_{0} x_{4} - x_{4} + x_{7}
\end{align*}
% END RESIDUAL 0
% BEGIN RESIDUAL 1
\begin{align*}
R_{1} ={}& -x_{2} - x_{4} x_{5} + x_{6} x_{7}
\end{align*}
% END RESIDUAL 1
% BEGIN RESIDUAL 2
\begin{align*}
R_{2} ={}& -1 + x_{1} - x_{6} x_{7}
\end{align*}
% END RESIDUAL 2
% BEGIN RESIDUAL 3
\begin{align*}
R_{3} ={}& -2x_{2} x_{213} - x_{3} x_{18} + x_{4} + 2x_{16}
\end{align*}
% END RESIDUAL 3
% BEGIN RESIDUAL 4
\begin{align*}
R_{4} ={}& -2 - x_{8} - x_{8}^{2} - 2x_{8} x_{9} - 2x_{9} - x_{9}^{2} \\
&{}  + x_{19}
\end{align*}
% END RESIDUAL 4
% BEGIN RESIDUAL 5
\begin{align*}
R_{5} ={}& -2 - x_{8} - x_{8}^{2} - 8x_{8} x_{19} - 4x_{8} x_{20} - 8x_{19} \\
&{}  - 16x_{19}^{2} - 16x_{19} x_{20} - 4x_{20} - 4x_{20}^{2} + x_{21}
\end{align*}
% END RESIDUAL 5
% BEGIN RESIDUAL 6
\begin{align*}
R_{6} ={}& 1 + 2x_{21} - 4x_{21}^{2} + x_{22}
\end{align*}
% END RESIDUAL 6
% BEGIN RESIDUAL 7
\begin{align*}
R_{7} ={}& -2 - x_{9}^{2} - 2x_{9} x_{10} - x_{10}^{2} + x_{23}
\end{align*}
% END RESIDUAL 7
% BEGIN RESIDUAL 8
\begin{align*}
R_{8} ={}& -1 + 2x_{8} - x_{8}^{2} - 2x_{8} x_{10} + 2x_{10} - x_{10}^{2} \\
&{}  + x_{24}
\end{align*}
% END RESIDUAL 8
% BEGIN RESIDUAL 9
\begin{align*}
R_{9} ={}& -7 + 33x_{8} - 35x_{8}^{2} - 96x_{8} x_{10} - 24x_{8} x_{19} - 4x_{8} x_{23} \\
&{}  - 24x_{8} x_{24} + 16x_{9} - 2x_{9}^{2} - 8x_{9} x_{23} + 32x_{10} - 64x_{10}^{2} \\
&{}  - 32x_{10} x_{19} - 32x_{10} x_{24} + 4x_{19} - 4x_{19}^{2} - 8x_{19} x_{24} + 12x_{23} \\
&{}  - 4x_{23}^{2} + 8x_{24} - 4x_{24}^{2} + x_{25}
\end{align*}
% END RESIDUAL 9
% BEGIN RESIDUAL 10
\begin{align*}
R_{10} ={}& -16 + 14x_{8} - 2x_{8}^{2} - 8x_{8} x_{9} - 8x_{8} x_{23} + 8x_{9} \\
&{}  + 8x_{23} - 4x_{25}^{2} + x_{26} - 2x_{43}
\end{align*}
% END RESIDUAL 10
% BEGIN RESIDUAL 11
\begin{align*}
R_{11} ={}& -6 + 16x_{10} - 16x_{10}^{2} - 16x_{10} x_{23} - 16x_{10} x_{24} + 8x_{23} \\
&{}  - 4x_{23}^{2} - 8x_{23} x_{24} + 8x_{24} - 4x_{24}^{2} + x_{27}
\end{align*}
% END RESIDUAL 11
% BEGIN RESIDUAL 12
\begin{align*}
R_{12} ={}& -5 + 24x_{10} - 16x_{10}^{2} - 16x_{10} x_{19} - 32x_{10} x_{24} - 16x_{10} x_{27} \\
&{}  + 10x_{19} - 4x_{19}^{2} - 16x_{19} x_{24} - 8x_{19} x_{27} + 24x_{24} - 16x_{24}^{2} \\
&{}  - 16x_{24} x_{27} + 12x_{27} - 4x_{27}^{2} + x_{28}
\end{align*}
% END RESIDUAL 12
% BEGIN RESIDUAL 13
\begin{align*}
R_{13} ={}& -12 - 29x_{8} - 37x_{8}^{2} - 24x_{8} x_{9} - 112x_{8} x_{10} - 4x_{8} x_{23} \\
&{}  - 24x_{8} x_{24} - 39x_{9} - 37x_{9}^{2} - 112x_{9} x_{10} - 24x_{9} x_{23} - 4x_{9} x_{24} \\
&{}  - 80x_{10} - 128x_{10}^{2} - 32x_{10} x_{23} - 32x_{10} x_{24} - 12x_{23} - 4x_{23}^{2} \\
&{}  - 8x_{24} - 4x_{24}^{2} + x_{29}
\end{align*}
% END RESIDUAL 13
% BEGIN RESIDUAL 14
\begin{align*}
R_{14} ={}& -18 - 10x_{8} - 2x_{8}^{2} - 24x_{8} x_{9} - 32x_{8} x_{10} - 8x_{8} x_{23} \\
&{}  - 72x_{9} - 72x_{9}^{2} - 192x_{9} x_{10} - 48x_{9} x_{23} - 96x_{10} - 128x_{10}^{2} \\
&{}  - 64x_{10} x_{23} - 24x_{23} - 8x_{23}^{2} - 4x_{29}^{2} + x_{30}
\end{align*}
% END RESIDUAL 14
% BEGIN RESIDUAL 15
\begin{align*}
R_{15} ={}& -32 - 42x_{8} + 8x_{8} x_{10} + 16x_{8} x_{23} - 84x_{9} + 16x_{9} x_{10} \\
&{}  + 32x_{9} x_{23} + 21x_{10} - 3x_{10}^{2} - 4x_{10} x_{19} - 8x_{10} x_{23} + 21x_{19} \\
&{}  - 8x_{19} x_{23} + 24x_{23} - 4x_{23}^{2} + x_{31} - x_{59}
\end{align*}
% END RESIDUAL 15
% BEGIN RESIDUAL 16
\begin{align*}
R_{16} ={}& -1 - 24x_{8} - 64x_{8}^{2} - 160x_{8} x_{10} - 64x_{8} x_{23} - 32x_{8} x_{24} \\
&{}  - 32x_{8} x_{31} - 26x_{10} - 100x_{10}^{2} - 80x_{10} x_{23} - 40x_{10} x_{24} - 40x_{10} x_{31} \\
&{}  - 8x_{23} - 16x_{23}^{2} - 16x_{23} x_{24} - 16x_{23} x_{31} - 6x_{24} - 4x_{24}^{2} \\
&{}  - 8x_{24} x_{31} - 4x_{31} - 4x_{31}^{2} + x_{32}
\end{align*}
% END RESIDUAL 16
% BEGIN RESIDUAL 17
\begin{align*}
R_{17} ={}& 1 - 4x_{19}^{2} + x_{20}
\end{align*}
% END RESIDUAL 17
% BEGIN RESIDUAL 18
\begin{align*}
R_{18} ={}& -2 - x_{8} - x_{8}^{2} + 8x_{8} x_{19} - 4x_{8} x_{20} + 8x_{19} \\
&{}  - 16x_{19}^{2} + 16x_{19} x_{20} - 4x_{20} - 4x_{20}^{2} + x_{33}
\end{align*}
% END RESIDUAL 18
% BEGIN RESIDUAL 19
\begin{align*}
R_{19} ={}& 1 + 2x_{33} - 4x_{33}^{2} + x_{34}
\end{align*}
% END RESIDUAL 19
% BEGIN RESIDUAL 20
\begin{align*}
R_{20} ={}& -x_{9}^{2} - 4x_{9} x_{34} - 4x_{34}^{2} + x_{35}
\end{align*}
% END RESIDUAL 20
% BEGIN RESIDUAL 21
\begin{align*}
R_{21} ={}& -8 - 20x_{8} - 16x_{8}^{2} - 64x_{8} x_{9} + 32x_{8} x_{19} - 40x_{9} \\
&{}  - 64x_{9}^{2} + 64x_{9} x_{19} + 18x_{19} - 16x_{19}^{2} + x_{36}
\end{align*}
% END RESIDUAL 21
% BEGIN RESIDUAL 22
\begin{align*}
R_{22} ={}& 8x_{9} - 62x_{9}^{2} - 92x_{9} x_{10} - 32x_{9} x_{23} - 32x_{9} x_{36} + 6x_{10} \\
&{}  - 34x_{10}^{2} - 24x_{10} x_{23} - 24x_{10} x_{36} - 4x_{23}^{2} - 8x_{23} x_{36} + 4x_{36} \\
&{}  - 4x_{36}^{2} + x_{37}
\end{align*}
% END RESIDUAL 22
% BEGIN RESIDUAL 23
\begin{align*}
R_{23} ={}& 8x_{8} - 64x_{8}^{2} - 96x_{8} x_{10} - 32x_{8} x_{24} - 32x_{8} x_{37} + 6x_{10} \\
&{}  - 36x_{10}^{2} - 24x_{10} x_{24} - 24x_{10} x_{37} + 2x_{24} - 4x_{24}^{2} - 8x_{24} x_{37} \\
&{}  + 4x_{37} - 4x_{37}^{2} + x_{38}
\end{align*}
% END RESIDUAL 23
% BEGIN RESIDUAL 24
\begin{align*}
R_{24} ={}& -x_{10}^{2} - 2x_{10} x_{11} - x_{11}^{2} + x_{39}
\end{align*}
% END RESIDUAL 24
% BEGIN RESIDUAL 25
\begin{align*}
R_{25} ={}& -32 + 22x_{8} - 3x_{8}^{2} - 8x_{8} x_{9} - 8x_{8} x_{23} - 8x_{8} x_{24} \\
&{}  - 4x_{8} x_{39} + 16x_{9} - 8x_{9} x_{24} - 2x_{10}^{2} - 8x_{10} x_{39} + 16x_{23} \\
&{}  - 8x_{23} x_{24} + 22x_{24} - 4x_{24}^{2} + 4x_{39} - 4x_{39}^{2} + x_{40} \\
&{}  - x_{43}
\end{align*}
% END RESIDUAL 25
% BEGIN RESIDUAL 26
\begin{align*}
R_{26} ={}& 6 - 16x_{8}^{2} - 36x_{8} x_{10} - 32x_{8} x_{24} - 40x_{8} x_{39} - 16x_{8} x_{40} \\
&{}  - 22x_{10}^{2} - 32x_{10} x_{24} - 48x_{10} x_{39} - 16x_{10} x_{40} - 2x_{24} - 16x_{24}^{2} \\
&{}  - 32x_{24} x_{39} - 16x_{24} x_{40} + 4x_{39} - 24x_{39}^{2} - 16x_{39} x_{40} - 4x_{40}^{2} \\
&{}  + x_{41}
\end{align*}
% END RESIDUAL 26
% BEGIN RESIDUAL 27
\begin{align*}
R_{27} ={}& x_{10} - 25x_{10}^{2} - 40x_{10} x_{24} - 20x_{10} x_{27} - 16x_{24}^{2} - 16x_{24} x_{27} \\
&{}  - 4x_{27}^{2} + x_{42}
\end{align*}
% END RESIDUAL 27
% BEGIN RESIDUAL 28
\begin{align*}
R_{28} ={}& -1 + 8x_{9} - 4x_{9}^{2} - 8x_{9} x_{23} + 8x_{23} - 4x_{23}^{2} \\
&{}  + x_{43}
\end{align*}
% END RESIDUAL 28
% BEGIN RESIDUAL 29
\begin{align*}
R_{29} ={}& -1 - 16x_{9}^{2} - 8x_{9} x_{10} - 32x_{9} x_{23} - 16x_{9} x_{43} - x_{10}^{2} \\
&{}  - 8x_{10} x_{23} - 4x_{10} x_{43} - 16x_{23}^{2} - 16x_{23} x_{43} - 4x_{43}^{2} + x_{44}
\end{align*}
% END RESIDUAL 29
% BEGIN RESIDUAL 30
\begin{align*}
R_{30} ={}& 3 + x_{9} - 2x_{9} x_{10} - 4x_{9} x_{44} - 3x_{10}^{2} - 8x_{10} x_{44} \\
&{}  - x_{23} - 4x_{44}^{2} + x_{45}
\end{align*}
% END RESIDUAL 30
% BEGIN RESIDUAL 31
\begin{align*}
R_{31} ={}& -1 - 4x_{45} - 4x_{45}^{2} + x_{46}
\end{align*}
% END RESIDUAL 31
% BEGIN RESIDUAL 32
\begin{align*}
R_{32} ={}& -1 - x_{10}^{2} - 8x_{10} x_{45} - 4x_{10} x_{46} - 16x_{45}^{2} - 16x_{45} x_{46} \\
&{}  - 4x_{46}^{2} + x_{47}
\end{align*}
% END RESIDUAL 32
% BEGIN RESIDUAL 33
\begin{align*}
R_{33} ={}& 2 + x_{9} - x_{9}^{2} - 4x_{9} x_{10} - 4x_{9} x_{47} - 4x_{10}^{2} \\
&{}  - 8x_{10} x_{47} - 4x_{47}^{2} + x_{48}
\end{align*}
% END RESIDUAL 33
% BEGIN RESIDUAL 34
\begin{align*}
R_{34} ={}& -4x_{42}^{2} - 8x_{42} x_{48} - 4x_{48}^{2} + x_{49}
\end{align*}
% END RESIDUAL 34
% BEGIN RESIDUAL 35
\begin{align*}
R_{35} ={}& -9 + 7x_{8} - x_{8}^{2} - 8x_{8} x_{42} - 4x_{8} x_{49} + 24x_{42} \\
&{}  - 16x_{42}^{2} - 16x_{42} x_{49} + 12x_{49} - 4x_{49}^{2} + x_{50}
\end{align*}
% END RESIDUAL 35
% BEGIN RESIDUAL 36
\begin{align*}
R_{36} ={}& -1 - 4x_{50} - 4x_{50}^{2} + x_{51}
\end{align*}
% END RESIDUAL 36
% BEGIN RESIDUAL 37
\begin{align*}
R_{37} ={}& -x_{10}^{2} - 8x_{10} x_{50} - 4x_{10} x_{51} - 16x_{50}^{2} - 16x_{50} x_{51} - 4x_{51}^{2} \\
&{}  + x_{52}
\end{align*}
% END RESIDUAL 37
% BEGIN RESIDUAL 38
\begin{align*}
R_{38} ={}& -x_{8} - x_{8}^{2} - 4x_{8} x_{10} - 4x_{8} x_{52} - 4x_{10} - 4x_{10}^{2} \\
&{}  - 8x_{10} x_{52} - 4x_{52} - 4x_{52}^{2} + x_{53}
\end{align*}
% END RESIDUAL 38
% BEGIN RESIDUAL 39
\begin{align*}
R_{39} ={}& -10 - 12x_{53} - 4x_{53}^{2} + x_{54}
\end{align*}
% END RESIDUAL 39
% BEGIN RESIDUAL 40
\begin{align*}
R_{40} ={}& -19 + 4x_{9} - 4x_{9}^{2} - 8x_{9} x_{23} - 16x_{9} x_{53} - 8x_{9} x_{54} \\
&{}  + 4x_{23} - 4x_{23}^{2} - 16x_{23} x_{53} - 8x_{23} x_{54} - 16x_{53} - 24x_{53}^{2} \\
&{}  - 16x_{53} x_{54} + 6x_{54} - 4x_{54}^{2} + x_{55}
\end{align*}
% END RESIDUAL 40
% BEGIN RESIDUAL 41
\begin{align*}
R_{41} ={}& -10 + x_{9} - 25x_{9}^{2} - 40x_{9} x_{23} - 20x_{9} x_{55} - 16x_{23}^{2} \\
&{}  - 16x_{23} x_{55} - 12x_{48} - 4x_{48}^{2} - 4x_{55}^{2} + x_{56}
\end{align*}
% END RESIDUAL 41
% BEGIN RESIDUAL 42
\begin{align*}
R_{42} ={}& 21 + 24x_{48} - 8x_{48}^{2} - 16x_{48} x_{56} - 2x_{56} - 4x_{56}^{2} \\
&{}  + x_{57}
\end{align*}
% END RESIDUAL 42
% BEGIN RESIDUAL 43
\begin{align*}
R_{43} ={}& -5 + 4x_{8} + 8x_{9} + x_{10} - x_{10}^{2} - 4x_{10} x_{57} \\
&{}  - 2x_{19} - 4x_{57}^{2} + x_{58} + 12x_{59} - 4x_{59}^{2}
\end{align*}
% END RESIDUAL 43
% BEGIN RESIDUAL 44
\begin{align*}
R_{44} ={}& -6x_{8} - 16x_{8}^{2} - 64x_{8} x_{9} + 8x_{8} x_{10} + 16x_{8} x_{19} - 12x_{9} \\
&{}  - 64x_{9}^{2} + 16x_{9} x_{10} + 32x_{9} x_{19} + 3x_{10} - x_{10}^{2} - 4x_{10} x_{19} \\
&{}  + 3x_{19} - 4x_{19}^{2} + x_{59}
\end{align*}
% END RESIDUAL 44
% BEGIN RESIDUAL 45
\begin{align*}
R_{45} ={}& -12 + 8x_{8} + 16x_{9} - 4x_{19} - 4x_{58}^{2} + 24x_{59} \\
&{}  - 8x_{59}^{2} + x_{60}
\end{align*}
% END RESIDUAL 45
% BEGIN RESIDUAL 46
\begin{align*}
R_{46} ={}& -1 - 4x_{60} - 4x_{60}^{2} + x_{61}
\end{align*}
% END RESIDUAL 46
% BEGIN RESIDUAL 47
\begin{align*}
R_{47} ={}& -x_{2} x_{13} + x_{62}
\end{align*}
% END RESIDUAL 47
% BEGIN RESIDUAL 48
\begin{align*}
R_{48} ={}& x_{12} - x_{13} - 2x_{22} x_{62} - x_{63}
\end{align*}
% END RESIDUAL 48
% BEGIN RESIDUAL 49
\begin{align*}
R_{49} ={}& x_{12} - x_{13} - 2x_{26} x_{62} - x_{64}
\end{align*}
% END RESIDUAL 49
% BEGIN RESIDUAL 50
\begin{align*}
R_{50} ={}& x_{12} - x_{13} - 2x_{28} x_{62} - 6x_{62} - x_{65}
\end{align*}
% END RESIDUAL 50
% BEGIN RESIDUAL 51
\begin{align*}
R_{51} ={}& x_{12} - x_{13} - 2x_{30} x_{62} - x_{66}
\end{align*}
% END RESIDUAL 51
% BEGIN RESIDUAL 52
\begin{align*}
R_{52} ={}& x_{12} - x_{13} - 2x_{32} x_{62} + 2x_{62} - x_{67}
\end{align*}
% END RESIDUAL 52
% BEGIN RESIDUAL 53
\begin{align*}
R_{53} ={}& x_{12} - x_{13} - 4x_{34} x_{62} - 2x_{35} x_{62} - x_{68}
\end{align*}
% END RESIDUAL 53
% BEGIN RESIDUAL 54
\begin{align*}
R_{54} ={}& x_{12} - x_{13} - 2x_{38} x_{62} - x_{69}
\end{align*}
% END RESIDUAL 54
% BEGIN RESIDUAL 55
\begin{align*}
R_{55} ={}& 2x_{28} x_{62} - 2x_{41} x_{62} - 6x_{62} + x_{65} - x_{70}
\end{align*}
% END RESIDUAL 55
% BEGIN RESIDUAL 56
\begin{align*}
R_{56} ={}& -x_{8} x_{9} + x_{8} x_{10} + x_{9} + x_{9} x_{10} - x_{10} - x_{10}^{2} \\
&{}  + x_{71}
\end{align*}
% END RESIDUAL 56
% BEGIN RESIDUAL 57
\begin{align*}
R_{57} ={}& 1 + 2x_{11} + x_{11}^{2} - x_{71}^{2} + x_{72}
\end{align*}
% END RESIDUAL 57
% BEGIN RESIDUAL 58
\begin{align*}
R_{58} ={}& -x_{12}^{2} + 2x_{12} x_{13} + 2x_{12} x_{74} - x_{13}^{2} - 2x_{13} x_{74} - x_{72}^{2} \\
&{}  + x_{73} - x_{74}^{2}
\end{align*}
% END RESIDUAL 58
% BEGIN RESIDUAL 59
\begin{align*}
R_{59} ={}& -4x_{60} x_{62} - 2x_{61} x_{62} + x_{74}
\end{align*}
% END RESIDUAL 59
% BEGIN RESIDUAL 60
\begin{align*}
R_{60} ={}& -x_{9} x_{62} - x_{13} + x_{75} x_{76} + x_{77}
\end{align*}
% END RESIDUAL 60
% BEGIN RESIDUAL 61
\begin{align*}
R_{61} ={}& -1 + 4x_{2} + 4x_{2} x_{8} + 8x_{2} x_{9} - 2x_{2} x_{19} + x_{75}
\end{align*}
% END RESIDUAL 61
% BEGIN RESIDUAL 62
\begin{align*}
R_{62} ={}& -x_{8} x_{208} - x_{12} + x_{76} + x_{208}
\end{align*}
% END RESIDUAL 62
% BEGIN RESIDUAL 63
\begin{align*}
R_{63} ={}& x_{9} x_{208} - 2x_{19} x_{62} + 2x_{32} x_{62} + x_{67} + x_{78}
\end{align*}
% END RESIDUAL 63
% BEGIN RESIDUAL 64
\begin{align*}
R_{64} ={}& -x_{12} x_{63} + x_{13} x_{63} + x_{79}
\end{align*}
% END RESIDUAL 64
% BEGIN RESIDUAL 65
\begin{align*}
R_{65} ={}& -x_{64} x_{79} + x_{80}
\end{align*}
% END RESIDUAL 65
% BEGIN RESIDUAL 66
\begin{align*}
R_{66} ={}& -x_{65} x_{80} + x_{81}
\end{align*}
% END RESIDUAL 66
% BEGIN RESIDUAL 67
\begin{align*}
R_{67} ={}& -x_{66} x_{81} + x_{82}
\end{align*}
% END RESIDUAL 67
% BEGIN RESIDUAL 68
\begin{align*}
R_{68} ={}& -x_{67} x_{82} + x_{83}
\end{align*}
% END RESIDUAL 68
% BEGIN RESIDUAL 69
\begin{align*}
R_{69} ={}& -x_{68} x_{83} + x_{84}
\end{align*}
% END RESIDUAL 69
% BEGIN RESIDUAL 70
\begin{align*}
R_{70} ={}& -x_{69} x_{84} + x_{85}
\end{align*}
% END RESIDUAL 70
% BEGIN RESIDUAL 71
\begin{align*}
R_{71} ={}& -x_{70} x_{85} + x_{86}
\end{align*}
% END RESIDUAL 71
% BEGIN RESIDUAL 72
\begin{align*}
R_{72} ={}& -x_{77} x_{86} + x_{87}
\end{align*}
% END RESIDUAL 72
% BEGIN RESIDUAL 73
\begin{align*}
R_{73} ={}& -x_{78} x_{87} + x_{88}
\end{align*}
% END RESIDUAL 73
% BEGIN RESIDUAL 74
\begin{align*}
R_{74} ={}& -x_{88}^{2} + x_{89}
\end{align*}
% END RESIDUAL 74
% BEGIN RESIDUAL 75
\begin{align*}
R_{75} ={}& -x_{4} - 2x_{4} x_{16} + 2x_{16} + x_{90}
\end{align*}
% END RESIDUAL 75
% BEGIN RESIDUAL 76
\begin{align*}
R_{76} ={}& -x_{14} x_{90} + x_{91}
\end{align*}
% END RESIDUAL 76
% BEGIN RESIDUAL 77
\begin{align*}
R_{77} ={}& -2x_{1}^{2} + 2x_{1} x_{12} + 2x_{1} x_{13} + 2x_{1} x_{91} + 2x_{1} x_{92} - x_{12}^{2} \\
&{}  - 2x_{12} x_{91} - x_{13}^{2} - 2x_{13} x_{92} + x_{17} x_{94} - x_{73} x_{89} - x_{91}^{2} \\
&{}  - x_{92}^{2}
\end{align*}
% END RESIDUAL 77
% BEGIN RESIDUAL 78
\begin{align*}
R_{78} ={}& -x_{4} x_{15} + x_{15} - x_{15} x_{90} + x_{92}
\end{align*}
% END RESIDUAL 78
% BEGIN RESIDUAL 79
\begin{align*}
R_{79} ={}& -x_{95} x_{98} + x_{98} + x_{100}
\end{align*}
% END RESIDUAL 79
% BEGIN RESIDUAL 80
\begin{align*}
R_{80} ={}& -x_{3} x_{121} + x_{96} x_{200} - x_{99} - x_{100} x_{121}
\end{align*}
% END RESIDUAL 80
% BEGIN RESIDUAL 81
\begin{align*}
R_{81} ={}& x_{101} x_{200} - x_{102} x_{124} - x_{103}
\end{align*}
% END RESIDUAL 81
% BEGIN RESIDUAL 82
\begin{align*}
R_{82} ={}& x_{104} x_{200} - x_{105} x_{124} - x_{106}
\end{align*}
% END RESIDUAL 82
% BEGIN RESIDUAL 83
\begin{align*}
R_{83} ={}& x_{107} x_{200} - x_{108} x_{124} - x_{109}
\end{align*}
% END RESIDUAL 83
% BEGIN RESIDUAL 84
\begin{align*}
R_{84} ={}& x_{110} x_{200} - x_{111} x_{124} - x_{112}
\end{align*}
% END RESIDUAL 84
% BEGIN RESIDUAL 85
\begin{align*}
R_{85} ={}& x_{113} x_{201} - x_{115} x_{124} - x_{116}
\end{align*}
% END RESIDUAL 85
% BEGIN RESIDUAL 86
\begin{align*}
R_{86} ={}& x_{117} x_{201} - x_{118} x_{124} - x_{119}
\end{align*}
% END RESIDUAL 86
% BEGIN RESIDUAL 87
\begin{align*}
R_{87} ={}& -x_{17} x_{121} + x_{124}
\end{align*}
% END RESIDUAL 87
% BEGIN RESIDUAL 88
\begin{align*}
R_{88} ={}& x_{120} - x_{122} x_{124} - x_{123}
\end{align*}
% END RESIDUAL 88
% BEGIN RESIDUAL 89
\begin{align*}
R_{89} ={}& -x_{124} x_{126} + x_{125} - x_{127}
\end{align*}
% END RESIDUAL 89
% BEGIN RESIDUAL 90
\begin{align*}
R_{90} ={}& -x_{124} x_{129} + x_{128} x_{200} - x_{130}
\end{align*}
% END RESIDUAL 90
% BEGIN RESIDUAL 91
\begin{align*}
R_{91} ={}& x_{131} - 2x_{131} x_{133} + x_{134}
\end{align*}
% END RESIDUAL 91
% BEGIN RESIDUAL 92
\begin{align*}
R_{92} ={}& -x_{131} x_{134} - x_{134} + x_{135}
\end{align*}
% END RESIDUAL 92
% BEGIN RESIDUAL 93
\begin{align*}
R_{93} ={}& -1 + x_{132}^{2} + x_{134}^{2} - x_{135}^{2}
\end{align*}
% END RESIDUAL 93
% BEGIN RESIDUAL 94
\begin{align*}
R_{94} ={}& 1 - x_{132} x_{137} - x_{136} + x_{138}
\end{align*}
% END RESIDUAL 94
% BEGIN RESIDUAL 95
\begin{align*}
R_{95} ={}& -1 + 2x_{136} - x_{136}^{2} + 2x_{136} x_{138} + x_{137}^{2}
\end{align*}
% END RESIDUAL 95
% BEGIN RESIDUAL 96
\begin{align*}
R_{96} ={}& 1 + x_{0}^{2} - 2x_{0} x_{132} + x_{139}
\end{align*}
% END RESIDUAL 96
% BEGIN RESIDUAL 97
\begin{align*}
R_{97} ={}& -x_{0} x_{182} + x_{2} + x_{139} x_{140} - x_{181}
\end{align*}
% END RESIDUAL 97
% BEGIN RESIDUAL 98
\begin{align*}
R_{98} ={}& 1 + x_{3}^{2} - 2x_{3} x_{132} + x_{141}
\end{align*}
% END RESIDUAL 98
% BEGIN RESIDUAL 99
\begin{align*}
R_{99} ={}& -x_{3} x_{173} + x_{17} + x_{141} x_{142} - x_{172}
\end{align*}
% END RESIDUAL 99
% BEGIN RESIDUAL 100
\begin{align*}
R_{100} ={}& -1 + 2x_{132} x_{184} - x_{143} - x_{184}^{2}
\end{align*}
% END RESIDUAL 100
% BEGIN RESIDUAL 101
\begin{align*}
R_{101} ={}& -1 + x_{95} + x_{143} x_{144} - x_{172} - x_{173} x_{184}
\end{align*}
% END RESIDUAL 101
% BEGIN RESIDUAL 102
\begin{align*}
R_{102} ={}& -x_{132} x_{146} - x_{145} + x_{147}
\end{align*}
% END RESIDUAL 102
% BEGIN RESIDUAL 103
\begin{align*}
R_{103} ={}& -1 - x_{145}^{2} + 2x_{145} x_{147} + x_{146}^{2}
\end{align*}
% END RESIDUAL 103
% BEGIN RESIDUAL 104
\begin{align*}
R_{104} ={}& 1 - x_{132} x_{149} - x_{148} + x_{150}
\end{align*}
% END RESIDUAL 104
% BEGIN RESIDUAL 105
\begin{align*}
R_{105} ={}& -1 + 2x_{148} - x_{148}^{2} + 2x_{148} x_{150} + x_{149}^{2}
\end{align*}
% END RESIDUAL 105
% BEGIN RESIDUAL 106
\begin{align*}
R_{106} ={}& -x_{132} x_{152} - x_{151} + x_{153}
\end{align*}
% END RESIDUAL 106
% BEGIN RESIDUAL 107
\begin{align*}
R_{107} ={}& -1 - x_{151}^{2} + 2x_{151} x_{153} + x_{152}^{2}
\end{align*}
% END RESIDUAL 107
% BEGIN RESIDUAL 108
\begin{align*}
R_{108} ={}& -x_{132} x_{155} - x_{154} + x_{156}
\end{align*}
% END RESIDUAL 108
% BEGIN RESIDUAL 109
\begin{align*}
R_{109} ={}& -1 - x_{154}^{2} + 2x_{154} x_{156} + x_{155}^{2}
\end{align*}
% END RESIDUAL 109
% BEGIN RESIDUAL 110
\begin{align*}
R_{110} ={}& -x_{132} x_{158} - x_{157} + x_{159}
\end{align*}
% END RESIDUAL 110
% BEGIN RESIDUAL 111
\begin{align*}
R_{111} ={}& -1 - x_{157}^{2} + 2x_{157} x_{159} + x_{158}^{2}
\end{align*}
% END RESIDUAL 111
% BEGIN RESIDUAL 112
\begin{align*}
R_{112} ={}& -x_{132} x_{161} - x_{160} + x_{162}
\end{align*}
% END RESIDUAL 112
% BEGIN RESIDUAL 113
\begin{align*}
R_{113} ={}& -1 - x_{160}^{2} + 2x_{160} x_{162} + x_{161}^{2}
\end{align*}
% END RESIDUAL 113
% BEGIN RESIDUAL 114
\begin{align*}
R_{114} ={}& 1 - x_{132} x_{164} - x_{163} + x_{165}
\end{align*}
% END RESIDUAL 114
% BEGIN RESIDUAL 115
\begin{align*}
R_{115} ={}& -1 + 2x_{163} - x_{163}^{2} + 2x_{163} x_{165} + x_{164}^{2}
\end{align*}
% END RESIDUAL 115
% BEGIN RESIDUAL 116
\begin{align*}
R_{116} ={}& -x_{132} x_{167} - x_{166} + x_{168}
\end{align*}
% END RESIDUAL 116
% BEGIN RESIDUAL 117
\begin{align*}
R_{117} ={}& -1 - x_{166}^{2} + 2x_{166} x_{168} + x_{167}^{2}
\end{align*}
% END RESIDUAL 117
% BEGIN RESIDUAL 118
\begin{align*}
R_{118} ={}& -x_{132} x_{170} - x_{169} + x_{171}
\end{align*}
% END RESIDUAL 118
% BEGIN RESIDUAL 119
\begin{align*}
R_{119} ={}& -1 - x_{169}^{2} + 2x_{169} x_{171} + x_{170}^{2}
\end{align*}
% END RESIDUAL 119
% BEGIN RESIDUAL 120
\begin{align*}
R_{120} ={}& 1 - x_{132} x_{173} - x_{172} + x_{174}
\end{align*}
% END RESIDUAL 120
% BEGIN RESIDUAL 121
\begin{align*}
R_{121} ={}& -1 + 2x_{172} - x_{172}^{2} + 2x_{172} x_{174} + x_{173}^{2}
\end{align*}
% END RESIDUAL 121
% BEGIN RESIDUAL 122
\begin{align*}
R_{122} ={}& -x_{132} x_{176} - x_{175} + x_{177}
\end{align*}
% END RESIDUAL 122
% BEGIN RESIDUAL 123
\begin{align*}
R_{123} ={}& -1 - x_{175}^{2} + 2x_{175} x_{177} + x_{176}^{2}
\end{align*}
% END RESIDUAL 123
% BEGIN RESIDUAL 124
\begin{align*}
R_{124} ={}& -x_{132} x_{179} - x_{178} + x_{180}
\end{align*}
% END RESIDUAL 124
% BEGIN RESIDUAL 125
\begin{align*}
R_{125} ={}& -1 - x_{178}^{2} + 2x_{178} x_{180} + x_{179}^{2}
\end{align*}
% END RESIDUAL 125
% BEGIN RESIDUAL 126
\begin{align*}
R_{126} ={}& 1 - x_{132} x_{182} - x_{181} + x_{183}
\end{align*}
% END RESIDUAL 126
% BEGIN RESIDUAL 127
\begin{align*}
R_{127} ={}& -1 + 2x_{181} - x_{181}^{2} + 2x_{181} x_{183} + x_{182}^{2}
\end{align*}
% END RESIDUAL 127
% BEGIN RESIDUAL 128
\begin{align*}
R_{128} ={}& -5 - 2x_{0} - 4x_{0}^{2} - 2x_{0} x_{2} - x_{1}^{2} - 2x_{2}^{2} \\
&{}  - x_{3}^{2} - x_{4}^{2} + 2x_{8} - x_{8}^{2} - x_{9}^{2} - x_{10}^{2} \\
&{}  - x_{11}^{2} - x_{12}^{2} - x_{13}^{2} - x_{14}^{2} - x_{15}^{2} - 4x_{16}^{2} \\
&{}  - x_{17}^{2} - x_{93}^{2} + 2x_{95} - 2x_{95}^{2} - x_{96}^{2} - x_{97}^{2} \\
&{}  - x_{101}^{2} - x_{104}^{2} - x_{107}^{2} - x_{110}^{2} - x_{113}^{2} - x_{114}^{2} \\
&{}  - x_{117}^{2} - x_{120}^{2} - x_{121}^{2} - x_{125}^{2} - x_{128}^{2} + x_{131} \\
&{}  - x_{184}^{2}
\end{align*}
% END RESIDUAL 128
% BEGIN RESIDUAL 129
\begin{align*}
R_{129} ={}& 2x_{3} - 4x_{3} x_{8} - 4x_{3} x_{9} - 4x_{3} x_{10} - 4x_{3} x_{11} - 4x_{3} x_{12} \\
&{}  - 4x_{3} x_{13} - 4x_{3} x_{14} - 4x_{3} x_{15} - 8x_{3} x_{16} + x_{184}
\end{align*}
% END RESIDUAL 129
% BEGIN RESIDUAL 130
\begin{align*}
R_{130} ={}& -128 + x_{0} - 4x_{185}^{2}
\end{align*}
% END RESIDUAL 130
% BEGIN RESIDUAL 131
\begin{align*}
R_{131} ={}& -x_{132} x_{187} + x_{196}
\end{align*}
% END RESIDUAL 131
% BEGIN RESIDUAL 132
\begin{align*}
R_{132} ={}& -1 + x_{186}^{2} + x_{187}^{2} - x_{196}^{2}
\end{align*}
% END RESIDUAL 132
% BEGIN RESIDUAL 133
\begin{align*}
R_{133} ={}& -x_{189} x_{190} + x_{197}
\end{align*}
% END RESIDUAL 133
% BEGIN RESIDUAL 134
\begin{align*}
R_{134} ={}& -1 + x_{188}^{2} + x_{189}^{2} - x_{197}^{2}
\end{align*}
% END RESIDUAL 134
% BEGIN RESIDUAL 135
\begin{align*}
R_{135} ={}& -x_{179}^{2} + x_{198}
\end{align*}
% END RESIDUAL 135
% BEGIN RESIDUAL 136
\begin{align*}
R_{136} ={}& x_{187} - 2x_{192} x_{198} - x_{198}
\end{align*}
% END RESIDUAL 136
% BEGIN RESIDUAL 137
\begin{align*}
R_{137} ={}& -1 - 4x_{187} - 8x_{187} x_{191} + x_{190}
\end{align*}
% END RESIDUAL 137
% BEGIN RESIDUAL 138
\begin{align*}
R_{138} ={}& -x_{132} - x_{186} x_{193} + x_{190}
\end{align*}
% END RESIDUAL 138
% BEGIN RESIDUAL 139
\begin{align*}
R_{139} ={}& -x_{180} - x_{186} x_{194} + x_{188}
\end{align*}
% END RESIDUAL 139
% BEGIN RESIDUAL 140
\begin{align*}
R_{140} ={}& -2x_{16} - 4x_{179} x_{195} + x_{189}
\end{align*}
% END RESIDUAL 140
% BEGIN RESIDUAL 141
\begin{align*}
R_{141} ={}& 9 + 2x_{0}^{2} + 2x_{0} x_{2} - 2x_{0} x_{137} - 2x_{0} x_{146} + x_{1}^{2} \\
&{}  - 2x_{1} x_{173} + x_{2}^{2} - 2x_{2} x_{146} + x_{4}^{2} - 2x_{4} x_{164} - 2x_{8} \\
&{}  + x_{8}^{2} - 2x_{8} x_{149} + x_{9}^{2} - 2x_{9} x_{152} + x_{10}^{2} - 2x_{10} x_{155} \\
&{}  + x_{11}^{2} - 2x_{11} x_{158} + x_{12}^{2} - 2x_{12} x_{161} + x_{13}^{2} - 2x_{13} x_{167} \\
&{}  + x_{14}^{2} - 2x_{14} x_{170} + x_{15}^{2} - 2x_{15} x_{176} - 4x_{16} x_{179} + x_{93}^{2} \\
&{}  - 2x_{93} x_{182} - x_{132} x_{199} + x_{138}^{2} - 2x_{147} + x_{147}^{2} + 2x_{149} \\
&{}  + x_{150}^{2} - 2x_{153} + x_{153}^{2} - 2x_{156} + x_{156}^{2} - 2x_{159} \\
&{}  + x_{159}^{2} - 2x_{162} + x_{162}^{2} + x_{165}^{2} - 2x_{168} + x_{168}^{2} \\
&{}  - 2x_{171} + x_{171}^{2} + x_{174}^{2} - 2x_{177} + x_{177}^{2} + 2x_{179}^{2} \\
&{}  - 2x_{180} + x_{180}^{2} + x_{183}^{2} + x_{199} - x_{202}^{2} - x_{203}^{2} \\
&{}  - x_{204} - x_{204}^{2}
\end{align*}
% END RESIDUAL 141
% BEGIN RESIDUAL 142
\begin{align*}
R_{142} ={}& -x_{97} x_{200} + x_{121}
\end{align*}
% END RESIDUAL 142
% BEGIN RESIDUAL 143
\begin{align*}
R_{143} ={}& -x_{114} x_{201} + x_{121}
\end{align*}
% END RESIDUAL 143
% BEGIN RESIDUAL 144
\begin{align*}
R_{144} ={}& -1 - 4x_{16}^{2} - x_{137}^{2} - x_{146}^{2} - x_{149}^{2} - x_{152}^{2} \\
&{}  - x_{155}^{2} - x_{158}^{2} - x_{161}^{2} - x_{164}^{2} - x_{167}^{2} - x_{170}^{2} \\
&{}  - x_{173}^{2} - x_{176}^{2} + x_{179}^{2} - x_{182}^{2} - x_{202}^{2} - x_{203}^{2} \\
&{}  - x_{204} - x_{204}^{2}
\end{align*}
% END RESIDUAL 144
% BEGIN RESIDUAL 145
\begin{align*}
R_{145} ={}& -1 - 4x_{99}^{2} - 4x_{103}^{2} - 4x_{106}^{2} - 4x_{109}^{2} - 4x_{112}^{2} \\
&{}  - 4x_{116}^{2} - 4x_{119}^{2} + x_{121}^{2} - 4x_{123}^{2} - 4x_{127}^{2} - 4x_{130}^{2} \\
&{}  - x_{205}^{2} - x_{206}^{2} - x_{207} - x_{207}^{2}
\end{align*}
% END RESIDUAL 145
% BEGIN RESIDUAL 146
\begin{align*}
R_{146} ={}& -x_{2} x_{12} + x_{208}
\end{align*}
% END RESIDUAL 146
% BEGIN RESIDUAL 147
\begin{align*}
R_{147} ={}& x_{95} x_{210} - 2x_{96} x_{146} - 2x_{101} x_{149} - 2x_{104} x_{152} - 2x_{107} x_{155} - 2x_{110} x_{158} \\
&{}  - 2x_{113} x_{161} - 2x_{117} x_{167} - 2x_{120} x_{170} - 2x_{125} x_{176} - 2x_{128} x_{179} - 2x_{132} x_{210} \\
&{}  + 2x_{146} x_{147} + 2x_{149} + 2x_{149} x_{150} + 2x_{152} x_{153} + 2x_{155} x_{156} + 2x_{158} x_{159} \\
&{}  + 2x_{161} x_{162} + 2x_{167} x_{168} + 2x_{170} x_{171} + 2x_{176} x_{177} + 2x_{179} x_{180} + x_{209}
\end{align*}
% END RESIDUAL 147
% BEGIN RESIDUAL 148
\begin{align*}
R_{148} ={}& -10 - x_{95} x_{209} + x_{96}^{2} - 2x_{96} x_{145} + x_{101}^{2} - 2x_{101} x_{148} \\
&{}  + x_{104}^{2} - 2x_{104} x_{151} + x_{107}^{2} - 2x_{107} x_{154} + x_{110}^{2} - 2x_{110} x_{157} \\
&{}  + x_{113}^{2} - 2x_{113} x_{160} + x_{117}^{2} - 2x_{117} x_{166} + x_{120}^{2} - 2x_{120} x_{169} \\
&{}  + x_{125}^{2} - 2x_{125} x_{175} + x_{128}^{2} - 2x_{128} x_{178} + 2x_{145} x_{147} + 2x_{148} \\
&{}  + 2x_{148} x_{150} + 2x_{151} x_{153} + 2x_{154} x_{156} + 2x_{157} x_{159} + 2x_{160} x_{162} + 2x_{166} x_{168} \\
&{}  + 2x_{169} x_{171} + 2x_{175} x_{177} + 2x_{178} x_{180} + x_{210}
\end{align*}
% END RESIDUAL 148
% BEGIN RESIDUAL 149
\begin{align*}
R_{149} ={}& x_{95} x_{212} - 2x_{97} x_{137} - 2x_{114} x_{164} - 2x_{121} x_{173} - 2x_{132} x_{212} + 2x_{137} \\
&{}  + 2x_{137} x_{138} + 2x_{164} + 2x_{164} x_{165} + 2x_{173} + 2x_{173} x_{174} + x_{211} \\
&{}  - x_{212}
\end{align*}
% END RESIDUAL 149
% BEGIN RESIDUAL 150
\begin{align*}
R_{150} ={}& -2 - x_{95} x_{211} + x_{97}^{2} - 2x_{97} x_{136} + x_{114}^{2} - 2x_{114} x_{163} \\
&{}  + x_{121}^{2} - 2x_{121} x_{172} + 2x_{136} + 2x_{136} x_{138} + 2x_{163} + 2x_{163} x_{165} \\
&{}  + x_{172}^{2} - x_{173}^{2} + x_{211} + x_{212}
\end{align*}
% END RESIDUAL 150
% BEGIN RESIDUAL 151
\begin{align*}
R_{151} ={}& -1 - x_{2} - 2x_{2} x_{213} + x_{4}
\end{align*}
% END RESIDUAL 151
% BEGIN RESIDUAL 152
\begin{align*}
R_{152} ={}& -x_{0}^{2} + x_{214}
\end{align*}
% END RESIDUAL 152
% BEGIN RESIDUAL 153
\begin{align*}
R_{153} ={}& x_{93} - x_{214}^{2}
\end{align*}
% END RESIDUAL 153
}


\section{Artifact map and reproducibility}

\begin{longtable}{@{}p{0.38\linewidth}p{0.57\linewidth}@{}}
\toprule Path & Purpose\\\midrule\endhead
\nolinkurl{report/report.tex}, PDF & Article, ordinary proofs and complete principal residuals.\\
\nolinkurl{equations/principal/} & Final quadratic residuals, collected polynomial, metrics, text and variable list.\\
\nolinkurl{equations/companions/} & Measured alternatives with smaller variable counts and different heights.\\
\nolinkurl{source/parent225/} & Pinned immediate parent and attribution.\\
\nolinkurl{source/anchor221/} & Complete new semantic source and literal change certificate.\\
\nolinkurl{stages/} & Exact projection and coordinate-change boundaries.\\
\nolinkurl{proofs/} & New unbounded proofs and independently assigned AI mathematical audits.\\
\nolinkurl{inherited/} & Identified inherited arithmetic, syntax and logical proof sources.\\
\nolinkurl{code/}, \nolinkurl{verify.py} & Deterministic producer, independent replay, evaluation and report checks.\\
\nolinkurl{verification/} & Receipts, exact source identities and finite checks actually run.\\
\nolinkurl{PROVENANCE.json}, \nolinkurl{SHA256SUMS} & Pinned origins and release integrity.\\
\bottomrule
\end{longtable}

Run \texttt{python3 -B verify.py} from the extracted package. The
verification and polynomial evaluation require only the Python standard
library. The PDF can be rebuilt with ordinary pdfLaTeX; timestamps and
TeX versions may change PDF bytes without changing any coefficient data.
See the README for the exact rebuild commands and recorded verification
scope. The provided evaluation program evaluates a specified tuple; it
does not decide whether a zero exists.

\begin{thebibliography}{99}
\bibitem{pooled} ProveIt, \emph{Pooling exact power certificates for
Diophantine consistency}, source 158, 25 October 2026.
\url{https://github.com/VladimirReshetnikov/ProveIt/tree/953fb4d82d4229f4efc5f9db4cc2bf38ae329bb4/Computability/DiophantineConsistency/con-zfc-diophantine-forms/158-pooled-powers}.
\bibitem{factorialfree} ProveIt, \emph{Removing factorial authenticity from
Diophantine consistency}, source 151, 25 October 2026.
\url{https://github.com/VladimirReshetnikov/ProveIt/tree/953fb4d82d4229f4efc5f9db4cc2bf38ae329bb4/Computability/DiophantineConsistency/con-zfc-diophantine-forms/151-factorial-free-binomial-crt}.
\bibitem{linearcrt} ProveIt, \emph{Linear CRT radii and semantic radix
recovery}, source 157, 25 October 2026. Pinned alongside the preceding source;
the copied \nolinkurl{MAIN_PROOF.md} supplies the sign bootstrap and
every-prime CRT completeness.
\bibitem{projection} ProveIt, \emph{Quadratic projection}, source 159,
25 October 2026. Exact common-Hessian and affine projection algorithms,
with an independent integer certificate checker.
\bibitem{anchor} ProveIt, \emph{Reusing a semantic Pell anchor over
unrestricted integers} and \emph{Recovering nonnegative arithmetic from
signed Pell orientation}, source 135, retained in source 158's provenance.
Copied as \nolinkurl{PROOF_315.md} and \nolinkurl{ORIENTATION_FIRST_PROOF.md}.
\bibitem{strictbudget} ProveIt, \emph{Aligning the amplification exponent
with a strict beta quotient range}, source 135, and source 158's
\emph{Cofinal strict bounds and a linear beta base}.
\bibitem{shell} ProveIt, \emph{Square-shell encoding proof} and the literal
\nolinkurl{shell_source.py} constructor source; identified by source 158's
\nolinkurl{SHELL_ENCODING_PROOF.md} supplement and exact source hashes.
\bibitem{calibration} ProveIt, Lean sources \nolinkurl{KCode222.lean},
\nolinkurl{KSeedDirect222.lean}, \nolinkurl{KEndpoint.lean}, and
\nolinkurl{ConstructibleChoice.lean}. Exact origins and inspected copies are
listed in \nolinkurl{PROVENANCE.json}; no new full Lean build is claimed.
\end{thebibliography}
\end{document}
